%%=============================================================================
%% Methodologie
%%=============================================================================

\chapter{\IfLanguageName{dutch}{Methodologie}{Methodology}}%
\label{ch:methodologie}

%% TODO: In dit hoofstuk geef je een korte toelichting over hoe je te werk bent
%% gegaan. Verdeel je onderzoek in grote fasen, en licht in elke fase toe wat
%% de doelstelling was, welke deliverables daar uit gekomen zijn, en welke
%% onderzoeksmethoden je daarbij toegepast hebt. Verantwoord waarom je
%% op deze manier te werk gegaan bent.
%% 
%% Voorbeelden van zulke fasen zijn: literatuurstudie, opstellen van een
%% requirements-analyse, opstellen long-list (bij vergelijkende studie),
%% selectie van geschikte tools (bij vergelijkende studie, "short-list"),
%% opzetten testopstelling/PoC, uitvoeren testen en verzamelen
%% van resultaten, analyse van resultaten, ...
%%
%% !!!!! LET OP !!!!!
%%
%% Het is uitdrukkelijk NIET de bedoeling dat je het grootste deel van de corpus
%% van je bachelorproef in dit hoofstuk verwerkt! Dit hoofdstuk is eerder een
%% kort overzicht van je plan van aanpak.
%%
%% Maak voor elke fase (behalve het literatuuronderzoek) een NIEUW HOOFDSTUK aan
%% en geef het een gepaste titel.

Zoals reeds in de inleiding vermeld, worden 3 luiken onderzocht/uitgevoerd. Daardoor zal deze bachelorproef verlopen in 3 fases.
\section{Fase 0 - Intern onderzoek}
Voor er begonnen kan worden aan het implementeren van de applicatie, moet er eerst goed begrepen worden op welke manier het probleem zich uit en wat eraan te doen is. Hiervoor wordt contact opgenomen met beide de productiechef en financieel bedienden binnen Freshfrozen. Zij bieden een goed antwoord op de deelvragen omtrent het probleemdomein doordat zij zelf zo dicht bij het probleem staan. Verder wordt gevraagd wat voor P. Kemseke (CEO van FreshFrozen) de 2 belangrijkste verbeteringen zouden zijn. Zo is er duidelijkheid over waar voor hem de focus moet liggen. Nadat al deze vragen gesteld zijn kan er pas echt van start gegaan worden.
\section{Fase 1 - Implementatie in het ERP-systeem}
In deze fase wordt er bekeken hoe het huidige ERP-systeem in elkaar zit. Hieronder valt onderzoek naar wat het systeem precies is en contact opnemen met het bedrijf die de software produceerde. Op die manier is er genoeg informatie om te weten of er met deze fase verder gewerkt kan worden. Het doel in deze fase is een proof-of-concept maken dat directe implementatie in een bestaand ERP-systeem mogelijk is. Ook wordt er gekeken naar de tijd dat dit zou kosten maar ook budget. De proof-of-concept bevat een duidelijk voorbeeld van een zelf geïmplementeerd systeem in de vorm van data die doorgegeven wordt binnen het systeem. In deze fase wordt er ook aangepast aan de tools die op dit moment gebruikt worden om het ERP-systeem te bouwen. Afhankelijk van de tools en de mogelijkheid tot directe implementatie kan deze fase langer of korter duren. In het geval van geen mogelijke implementatie zal deze fase lopen over een termijn van 5 dagen. In het geval van mogelijke implementatie zal deze fase lopen over een termijn van 10 dagen.
\section{Fase 2 - Communicatie met het ERP-Systeem}
Hier wordt onderzocht of communicatie met het ERP-Systeem een optie is. Dit door middel van API-calls die van en naar het systeem worden gestuurd. Ook hier wordt gekeken naar tijd en budget. Het doel van deze fase is om opnieuw een proof-of-concept te maken. Deze keer bevat deze een duidelijk voorbeeld van communicatie tussen een eigen uitwerking en het ERP-systeem. Zo kan het ERP-systeem API-calls gebruiken om op een externe database CRUD operations te doen. Ook de mogelijkheid om dit te doen moet onderzocht en besproken worden met de leverancier van het ERP-systeem. Voor het uitwerken van deze fase wordt gebruik gemaakt van een .NET backend die werkt als basis voor de CRUD operations. Dit wordt ook het aanspreekpunt van het ERP-systeem om data op te halen door middel van API-calls. Deze fase loopt over een termijn van 10 dagen.
\section{Fase 3 - Standalone applicatie}
Na het uitwerken van de vorige 2 fases volgt de implementatie van een standalone uitwerking die het probleem van Freshfrozen oplost. Het budget wordt hierbij opnieuw onderzocht. Het doel van deze fase is een standalone applicatie uitwerken die aan alle criteria van de onderzoeksvraag voldoet. Deze zal gebruik maken van de industriestandaard die terug te vinden is in de literatuurstudie. Deze fase zal lopen over een termijn van 15 dagen.

